\chapter{运营与推广策略}

\section{整体运营策略}

针对小农户「分散、下沉、信任半径短、熟人社会」的特征，本项目采取\hlgreen{「线下渠道打样板、线上内容做放大、政企联动做背书」}的组合运营策略，以「县域样板——跨县复制——区域连片」的节奏推进，避免盲目撒网。

\section{用户获取策略}

\subsection{线下渠道：农资终端与村镇基层合作}

农资店是农户购买农资的必经节点，也是农户最信任的「技术入口」。本项目与县域农资经销商、零售终端建立深度合作，在门店放置推广物料、由店主引导首单体验，并以「门店业绩分成」绑定利益，实现低成本、高信任的获客。

\subsection{线上渠道：下沉市场内容传播}

面向快手、抖音、微信视频号等下沉市场主阵地，生产「一张照片识别病虫害」的短视频与实操内容，以「实用、接地气、能救命」的内容调性实现自然裂变；在农业类微信群、乡镇社群中开展「晒病害、得方案」互动，形成传播闭环。

\subsection{政企联动：基层农技体系推广}

联合地方农业农村部门、植保站、供销社，将本产品纳入「数字乡村」「统防统治」示范项目，借助官方渠道批量触达农户，同时为后续政府项目采购埋下商业线索。

\section{用户留存与活跃运营}

\subsection{预警服务粘性运营}

病虫害预警具有「当季、当区、高频」特征，是天然的留存抓手。通过持续的本地化预警推送，让用户形成「每到关键农时就看慧农识病」的依赖，把低频的「看病」工具升级为高频的「防病」参谋。

\subsection{社群答疑与农户互助运营}

建立县域农户互助社群，由农技专家定期答疑，鼓励农户之间分享防治经验与成功案例，以「互助＋专家」双引擎提升活跃度与信任度，沉淀私域用户资产。

\section{产业合作生态建设}

纵向整合农资生产商、经销商、农技专家、基层农技体系，横向拓展气象、遥感、农业保险等数据合作伙伴，逐步构建「\hlgreen{诊断—用药—防治—保险}」的产业服务生态，使本项目从单一工具进化为区域农业植保服务的数字化基础设施。

\begin{figure}[H]
  \centering
  \includegraphics[width=0.86\linewidth]{user_growth.png}
  \caption{用户增长预测（累计注册用户与月活跃用户，万人）}
\end{figure}
